Requirements for Vision Section (550 words maximum)
Explain how your proposed work:
is of excellent quality and importance within or beyond the field(s) or area(s)
has the potential to advance current understanding, or generate new knowledge, thinking or discovery within or beyond the field or area of its focus
is timely, given current trends, context, and needs
impacts world-leading research, society, the economy or the environment 

Within the Vision section we also expect you to:
identify the potential direct or indirect benefits and who the beneficiaries might be
identify potential improvements in human or population health, whether through contributing to relieving disease or disability burden, improving quality of life or providing benefit to the health service or health-related industry
outline your plans for engagement, communication and dissemination about your research and its outcomes with the research community and, where appropriate, with potentially interested wider audiences


---------------------------

Vision 

Short summary of vision
Infection and microbiome interactions at the respiratory mucosa drive both health and respiratory disease, but we lack scalable tools to connect pathogen genotype to AMR, virulence and clinical outcomes at the level needed for personalised care.  World leading bioinformatics currently does not translate easily into new patient care with testing stuck in the 1900's.  
This professional fellowship, embedded in CCLI and the Floto Laboratory at University of Cambridge, is designed to systematically address the current limitations to extending world leading basic bioinformatics into scalable predictions for AMR, virulence and clinical outcomes from genomic data using machine learning.  

Excellent Quality
The Floto Laboratory, CCLI and collaborators at the Sanger at world leaders in understanding and predicting behaviour of globally distributed pathogens.  

Publications on pathoadaption, saltatory leaps and spread of multi-drug resistant bacteria. References to Science papers (Bryant, Bryant, Weimann, Boeck). Understanding limitations beyond AMR of failure of treatment and establishment of chronic infections / flairs (intracellular pathogens, disabling macrophage activity, changes in transcription and regulation) which lead to antimicrobial failures, clinical failures, deaths and spread of multi-resistant and untreatable diseases (references).

Bacformer - unprecedented capacity to predict, won numerous awards and currently in review for publication.  My own role at intersection of clinical medicine, bioinformatics and machine learning. As second author on our award winning work, Bacformer—developing AMR prediction and gene discovery—and principal investigator of the SputaGen prospective clinical study, I bring together direct clinical experience in managing patients on intravenous antibiotics, advanced training in mathematics and machine learning, and authorship across all underpinning technologies, positioning this fellowship to translate these discoveries into validated clinical tools (reference).

Importance and Generation of New Knowledge
Mucosal immunity and host–pathogen interactions lie at a critical junction in modern medicine. The respiratory epithelium is the primary interface between pathogen and host across the full spectrum of acute and chronic lung disease—community-acquired and nosocomial pneumonia, epidemic and pandemic respiratory infection, asthma, ARDS, bronchiectasis, cystic fibrosis, and interstitial lung disease. 

Emphasize the scale of the problem in dollars and patient deaths - community and hospital. 

Despite this breadth, clinical practice relies on binary susceptibility testing that routinely fails to predict treatment response. The fundamental scientific gap is our inability to predict pathogen behaviour—resistance, tolerance, virulence potential, and within-host evolutionary trajectory—or to identify the genomic and regulatory drivers that determine it. 

Timeliness
Clinically timely - Heightened global awareness of respiratory pandemic vulnerability, combined with the escalating epidemic of AMR and untreatable infections—both WHO-designated priority threats—creates a compelling imperative for this work now. 

This fellowship builds on a sequence of landmark publications from our group that together provide a uniquely strong foundation.

Knowledge of HGT, pathoadaption, hot-spot evolution. 

Genomic technology.

Machine learning - Bacformer.  

Bryant et al. (Science, 2021) demonstrated that M. abscessus undergoes stepwise pathogenic evolution via horizontally acquired epigenetic regulators and convergent mutation in a constrained pathoadaptive gene network, establishing a generalisable genome-based framework for the emergence of virulence from environmental organisms. Weimann et al. (Science, 2024) extended this to P. aeruginosa: analysis of 9,829 isolates identified 21 global epidemic clones generated by HGT events in transcriptional regulators, 224 pathoadaptive genes, and DksA1-regulated host specificity mediating CF versus non-CF infection preference—providing a blueprint for decoding genome-encoded host adaptation. Boeck et al. (Nature Microbiology, 2026) then demonstrated that antibiotic killing dynamics and tolerance are genetically encoded, heritable traits that predict individual patient outcomes independently of conventional resistance metrics, establishing a critical and previously underexplored dimension of bacterial susceptibility. 

Underpinning the entire analytical framework, Wiatrak, Abelson et al. (Science, submitted) developed Bacformer—a genome-scale foundation model trained on over 1.3 million bacterial genomes—which learns the functional syntax of bacterial evolution, enables end-to-end genotype-to-phenotype prediction with gene-level attribution, and achieves near-perfect AMR prediction on held-out isolates. 

We also need to explain the new availability of large public databases and metadata including on AMR phenotypes (Klebsiella and EBI) and 

The maturation of SputaGen as a set provides clinical data and metadata as well as fully collected sputum sample biobank to test hypotheses and challenge next steps of development in real clinical samples requirement assemblies and metagenomics. 


Impact and Beneficiaries
Outcome - A model that can predict AMR, virulence and assessment of utility in clinical predictions, ready for next steps of translation into a clinical tool with organised data and pipeline to test from benchtop to clinical diagnostics.  

Downstream from this, the hope is that patients with acute and chronic lung disease are the target beneficiaries, through individually targeted antibiotic selection, reduced treatment failure, and lower disease burden across WHO-priority conditions. 

CCLI, The Floto-lab and my position and career in this will be ready. 
Engagement, Communication, and Dissemination
Findings will be disseminated through high-impact peer-reviewed journals and international conferences with open-source release of all genomes, curated and prepared metadata, test pipelines, models and code as a resource for the scientific community to tackle this fundamental problem in an organised and systematic way. 

We will leverage broad public and professional interest in WHO AMR priority topics through traditional news outlets, targeted social media and institutional communications. 

This will advance the data set, testing pipeline and models to engage academic partners and industry, NHS clinical networks and commercial industry partners in this WHO-priority area—establishing the UK as a global leader in AI-driven infection prediction, personalised antibiotic therapy, and AMR diagnostics.

References
WHO AMR, Mucosal immunity, Boeck and AST / antibiotic killing, Byrant x 2, Weimann x 1, Boeck lab x 1, Wiatrak (in submission)
 (https://www.who.int/news-room/fact-sheets/detail/antimicrobial-resistance).
